\section{A Framework for Testing Structural Attribution}
\label{sec:method}

\subsection{A common form for score-weighted updates}
\label{sec:form}

Let $V=\{1,\dots,N\}$ index a population of candidate solutions with positions
$X_t \in \mathbb{R}^{N \times D}$, whose $i$th row $x_{i,t}$ is one solution,
and let $G_t=(V,E_t)$ be the interaction graph at iteration $t$ with
neighbourhoods $\mathcal{N}_i$ and degrees $d_i=|\mathcal{N}_i|$.

A broad family of population methods---cellular evolutionary
algorithms~\cite{alba2005exploration}, neighbourhood-topology particle
swarms~\cite{kennedy2002population}, localized consensus
dynamics~\cite{bungert2024polarized}, and the centrality-weighted optimizer we
study~\cite{jalali2026social}---share the update
%
\begin{equation}
\label{eq:master}
X_{t+1} \;=\; (1-\eta_t)\,X_t \;+\; \eta_t\Bigl[\, \textstyle\sum_{r} \theta_{r,t}\, W^{(r)}_t X_t \;+\; B_t \,\Bigr] \;+\; \Xi_t ,
\end{equation}
%
where each $W^{(r)}_t \in \mathbb{R}^{N\times N}$ is row-stochastic and
supported on $E_t$, the $\theta_{r,t}$ are non-negative schedule weights,
$B_t$ collects rank-one drift terms toward incumbents, and $\Xi_t$ is a
mutation or noise term. Everything specific to a score-weighted method sits
in how the mixing operators $W^{(r)}_t$ are built.

The construction under test assigns each node a structural score
$c_t : V \to \mathbb{R}_{\ge 0}$---betweenness centrality in our primary case
study---and weights neighbours in proportion to it:
%
\begin{equation}
\label{eq:W}
\bigl[W(c)\bigr]_{ij} \;=\;
\frac{\bigl(c(j)+\varepsilon\bigr)\,\mathbf{1}\{j \in \mathcal{N}_i\}}
     {\sum_{k \in \mathcal{N}_i}\bigl(c(k)+\varepsilon\bigr)} ,
\end{equation}
%
with $\varepsilon>0$ guarding against neighbourhoods of all-zero score. Write
$w_i = \bigl(W(c)\bigr)_{i,\cdot}$ restricted to $\mathcal{N}_i$, and let
$u_i = d_i^{-1}\mathbf{1}$ denote the uniform weighting, i.e.\ the operator
$W_{\mathrm{unif}}$ of the same graph with no structural score at all.

\subsection{What the structural claim asserts, and what it does not}
\label{sec:claim}

The published justification for \eqref{eq:W} is that betweenness identifies
agents positioned to bridge regions of the search space, so that learning
preferentially from them propagates useful information. Formally this is a
claim about the \emph{assignment} $v \mapsto c(v)$: that the pairing of scores
to graph positions, and thereby to the solutions those positions hold, is
informative about where to search.

Two distinct properties of \eqref{eq:W} could each produce a performance
difference, and they are routinely conflated:
%
\begin{enumerate}
\item[(P1)] \textbf{Heterogeneity.} The weights $w_i$ are unequal, so the
  update is not the neighbourhood mean. This changes the operator's spectrum
  and the effective mixing rate regardless of which node gets which score.
\item[(P2)] \textbf{Assignment.} \emph{Which} neighbour receives the large
  weight correlates with search-relevant structure.
\end{enumerate}
%
The claim being made is (P2). A conventional ablation---replacing $W(c)$ by
$W_{\mathrm{unif}}$---removes (P1) and (P2) simultaneously and therefore
cannot adjudicate it. This is the gap our controls close.

\subsection{Control 1: value permutation}
\label{sec:perm}

Let $\pi$ be a permutation of $V$ drawn uniformly at random, independently of
the search state, and resampled whenever the graph changes. The permutation
control replaces $W(c)$ by $W(c \circ \pi)$, leaving \eqref{eq:master},
the graph, the schedule, the budget and the random stream untouched.

The multiset of scores $\{\!\{c(v)\}\!\}_{v \in V}$ is preserved exactly, as is
the arithmetic cost of forming the operator. Destroyed is precisely the
correspondence (P2). The following two facts make this the right control.

\begin{proposition}[Uniform first moment]
\label{prop:mean}
For every $i$ and every $j \in \mathcal{N}_i$,
$\;\mathbb{E}_\pi\bigl[W(c \circ \pi)\bigr]_{ij} = 1/d_i$; equivalently
$\mathbb{E}_\pi\bigl[W(c \circ \pi)\bigr] = W_{\mathrm{unif}}$.
\end{proposition}

\begin{IEEEproof}
Fix $i$. Under a uniform permutation the scores
$\bigl(c(\pi(j))\bigr)_{j \in \mathcal{N}_i}$ form an exchangeable random
vector, being a sample without replacement from
$\{\!\{c(v)\}\!\}_{v\in V}$. The map in \eqref{eq:W} is a ratio of a
coordinate to the sum of coordinates and hence equivariant under permutation
of $\mathcal{N}_i$, so $\bigl(w_{ij}\bigr)_{j\in\mathcal{N}_i}$ is itself
exchangeable. Exchangeable coordinates have equal means, and they sum to one
almost surely; therefore each mean equals $1/d_i$.
\end{IEEEproof}

Proposition~\ref{prop:mean} says the control sits exactly between the two
alternatives usually compared: \emph{in expectation} it is the uniform
operator, yet \emph{in every realisation} it retains the weight heterogeneity
of the original. It therefore holds (P1) fixed in distribution while removing
(P2). A difference between $W(c)$ and $W(c\circ\pi)$ is attributable to the
assignment and to nothing else.

\subsection{When is the mechanism strong enough to matter?}
\label{sec:leverage}

A control is only informative if the mechanism it neutralises was doing
something. We quantify this directly. Let
$\bar{x}^{\,\mathrm{cent}}_i = \sum_{j} w_{ij} x_j$ and
$\bar{x}^{\,\mathrm{unif}}_i = d_i^{-1}\sum_{j} x_j$, and define the
per-node displacement $\Delta_i = \bar{x}^{\,\mathrm{cent}}_i -
\bar{x}^{\,\mathrm{unif}}_i$.

\begin{proposition}[Leverage bound]
\label{prop:leverage}
Let $s_i^2 = d_i^{-1}\sum_{j\in\mathcal{N}_i}\|x_j-\bar{x}^{\,\mathrm{unif}}_i\|^2$
be the neighbourhood dispersion and define the \emph{structural leverage}
%
\begin{equation}
\label{eq:leverage}
\chi_i \;=\; \sqrt{d_i}\,\bigl\|w_i-u_i\bigr\|_2 \;\in\; \bigl[0,\;\sqrt{d_i-1}\,\bigr].
\end{equation}
%
Then $\|\Delta_i\| \le \chi_i\, s_i$, with $\chi_i=0$ if and only if
$w_i=u_i$ and $\chi_i=\sqrt{d_i-1}$ if and only if $w_i$ is a point mass.
\end{proposition}

\begin{IEEEproof}
Since $\sum_{j\in\mathcal{N}_i}(w_{ij}-d_i^{-1})=0$, we may centre at
$\bar{x}^{\,\mathrm{unif}}_i$ and write
$\Delta_i=\sum_{j}(w_{ij}-d_i^{-1})\bigl(x_j-\bar{x}^{\,\mathrm{unif}}_i\bigr)$.
The Cauchy--Schwarz inequality for vector-valued sums gives
$\|\Delta_i\| \le \|w_i-u_i\|_2\bigl(\sum_j \|x_j-\bar{x}^{\,\mathrm{unif}}_i\|^2\bigr)^{1/2}
= \|w_i-u_i\|_2\sqrt{d_i}\,s_i$. For the extremes, $\|w_i-u_i\|_2=0$ iff
$w_i=u_i$, and over the simplex $\|w_i-u_i\|_2$ is maximised at a vertex,
where $\|w_i-u_i\|_2^2 = 1-d_i^{-1}$.
\end{IEEEproof}

Equation~\eqref{eq:leverage} separates the displacement into a purely
structural factor $\chi_i$, determined by how concentrated the weights are,
and a purely populational factor $s_i$. Alongside every control comparison we
report the normalised displacement
%
\begin{equation}
\label{eq:relshift}
\rho \;=\; \frac{1}{N}\sum_{i \in V} \frac{\|\Delta_i\|}{\sigma_{\mathrm{pop}}},
\qquad
\sigma_{\mathrm{pop}}^2 \;=\; \frac{1}{N}\sum_{i\in V}\bigl\|x_i-\bar{x}\bigr\|^2 ,
\end{equation}
%
which measures the mechanism's effect on the update in units of the
population's own spread. The interesting regime, and the one we in fact
observe, is $\rho$ large with the control nevertheless indistinguishable: a
mechanism that moves the search substantially while carrying no usable
information.

Two structural remarks follow from \eqref{eq:leverage} and bear on the
implementation. First, $\chi_i$ is large precisely when a few neighbours
dominate; because betweenness on sparse small-world graphs is exactly zero at
many nodes, the guard $\varepsilon$ in \eqref{eq:W} assigns those neighbours
weight $O(\varepsilon)$ and the operator behaves less like soft weighting than
like hard masking. Second, $\chi_i$ vanishes identically on vertex-transitive
graphs---the ring lattice being the degenerate case---so any such topology
renders the mechanism inert by construction, independently of the landscape.

\subsection{A positive control}
\label{sec:positive}

A null result invites the objection that the test cannot detect anything. The
objection is answered by applying the same control to a score that is
informative by construction.

The optimizer under study weights neighbours by two scores, not one. Besides
centrality it forms a second aggregate $\bar{x}^{\,\mathrm{inf}}$ using a
fitness-influence score $\phi(v)$, monotone in rank, so that fitter neighbours
carry more weight. That score \emph{is} search-relevant: it is a direct
function of the objective values the run has already paid for. Permuting
$\phi$ across nodes---identical machinery, identical value multiset, identical
cost---should therefore cost performance.

This gives a control pair with opposite predictions under the same test: a
permutation that must hurt if the instrument works at all, and a permutation
whose effect is the open question. Reporting both together separates ``the
mechanism carries no information'' from ``our test detects nothing.''

\subsection{Control 2: policy randomisation}
\label{sec:rand}

Adaptive components select among alternatives $\mathcal{A}=\{a_1,\dots,a_m\}$
using feedback accumulated during the run: an operator, a neighbourhood size,
a parameter value. Let $\Psi_t : \mathcal{H}_t \to \mathcal{A}$ be the
selection rule, a function of the observed history $\mathcal{H}_t$. The claim
attached to such a component is that $\Psi_t$ is \emph{informed}, i.e.\ that
conditioning on $\mathcal{H}_t$ improves the choice.

The policy control replaces $\Psi_t$ by $\Psi^{\mathrm{rand}}_t \sim
\mathrm{Unif}(\mathcal{A})$, holding the alternatives, the switching schedule
and the switching frequency fixed. It is the exact analogue of
Proposition~\ref{prop:mean} at the level of policies: the marginal
distribution over which alternatives are used is flattened, but the fact of
alternating among them is preserved. An adaptive rule that does not beat
$\Psi^{\mathrm{rand}}$ is not adapting, and whatever benefit the component
confers belongs to variation per se rather than to selection.

\subsection{Statistical protocol}
\label{sec:stats}

Because the claim under test is an equivalence, a non-significant hypothesis
test is not by itself evidence: it is equally consistent with an underpowered
study. We therefore report effect sizes with interval estimates as the primary
quantity, add two one-sided
tests~\cite{schuirmann1987comparison,lakens2017equivalence}, and treat
difference tests as secondary.

Runs are paired by seed: a variant and its control receive the same initial
population and the same random stream. Per function we report the median
outcome, Cliff's $\delta_f$~\cite{cliff1993dominance} with a percentile
bootstrap interval, and a two-sided Wilcoxon signed-rank test under Holm's
step-down correction over the functions of a suite~\cite{holm1979simple},
following established practice for this
setting~\cite{derrac2011practical,garcia2009study}. When more than two variants
are compared we apply a Friedman test per function and report mean ranks.

Two choices in that protocol deserve stating explicitly.

\emph{The effect size and the pairing.} Cliff's $\delta$ estimates stochastic
dominance, $P(X{>}Y)-P(X{<}Y)$, between the two outcome \emph{distributions};
it is computed over all cross-pairs and is therefore an unpaired statistic
even though the runs are paired. This is deliberate. The substantive question
is whether one arm tends to return worse results than the other, not whether it
does so within a shared random stream, and the unpaired form is invariant to
any monotone rescaling of the objective---indispensable on a suite whose errors
span eighteen orders of magnitude, where a paired mean difference is
uninterpretable. Pairing is used where it adds power and validity: the Wilcoxon
test operates on within-seed differences, and the bootstrap resamples seed
pairs rather than runs, so the intervals inherit the paired design. A reader
who prefers a strictly paired effect size can recover one from the per-run
records, which are released.

\emph{The equivalence margin.} We adopt $|\delta|<0.147$ for a negligible
effect and $|\delta|<0.33$ for a small one, the conventional thresholds for
this statistic~\cite{romano2006exploring}, rather than choosing a margin to
suit the data. Because that choice is a judgement and not a measurement, every
equivalence claim is also reported against a range of margins
(Section~\ref{sec:res:tost}), so a reader who prefers a stricter or looser
bound can read the conclusion off directly. No claim in this paper depends on
the margin taking one particular value.
